\section{Introduction}
\label{sec:intro}

The modern deep learning ecosystem is increasingly decentralized. Rather than training a monolith from scratch for every multi-task requirement, practitioners frequently leverage Parameter-Efficient Fine-Tuning (PEFT) techniques, such as Low-Rank Adaptation (LoRA) \cite{Hu2021LoRA}, to specialize pre-trained backbone models on various target domains. Consequently, a vibrant field of study has emerged around \emph{model merging} \cite{Choshen2022Merging,Wortsman2022ModelSoups}, which seeks to average, blend, or align these independently trained weights into a unified multi-task model without further training compute. Successful approaches range from naive weight averaging to sophisticated interference-reduction schemes such as Task Arithmetic \cite{Ilharco2022TaskArithmetic}, TIES-Merging \cite{Yadav2023TIES}, and DARE \cite{Yu2023DARE}. 

However, in real-world deployment pipelines, model merging does not happen in a vacuum. To deploy these multi-task models onto resource-constrained edge devices, they must undergo compression, typically via Post-Training Quantization (PTQ) \cite{Dettmers2022LLMint8,Frantar2022GPTQ,Lin2023AWQ} down to 8-bit or 4-bit precision. Under standard workflows, adapters trained via high-precision or dequantized pipelines (such as QLoRA \cite{Dettmers2023QLoRA}) are merged into the base model in high-precision (FP16/FP32), and the resulting fused model is re-quantized back to low-bit constraints for inference.

In this work, we argue that the deep learning community has succumbed to a dangerous methodological blindspot: **the full-precision model merging abstraction**. While model merging papers routinely report impressive gains in full-precision, they rarely evaluate their merged models under the post-training quantization constraints that are mandatory for actual deployment. We investigate the potential consequence of this blindspot: the \textbf{``Re-Quantization Silence''}, where post-hoc re-quantization completely obliterates the subtle task-specific updates encoded in low-rank adapters, reducing task-specific performance back to the level of the unadapted base model.

\begin{figure}[t]
  \centering
  \includegraphics[width=0.48\textwidth]{quantization_comparison.png}
  \caption{\textbf{Multi-Task Mean Accuracy across Quantization Configurations.} While Naive Re-Quantization (Naive-RQ) is virtually lossless under standard per-channel configurations, it suffers from a catastrophic performance drop under aggressive per-tensor 4-bit constraints. Optimization-based approaches (AdaMerging and QA-ACS) are highly robust and achieve the best overall accuracy, whereas SAWS provides a zero-data alternative but can degrade performance under per-tensor grids.}
  \label{fig:teaser}
\end{figure}

The mathematical core of this failure lies in a severe magnitude mismatch. The pre-trained base weight matrix $W_0$ typically exhibits a wide dynamic range, whereas the low-rank update matrix $\Delta W_{\text{merged}}$ contains fine-grained, low-magnitude task-specific updates. When these weights are added in high-precision, the combined weight matrix is dominated by $W_0$. When the merged weight matrix is subsequently passed through a low-bit uniform quantizer, the step size (scale factor $s$) is computed based on the maximum magnitude of the fused matrix (which is almost identical to that of $W_0$). Consequently, the tiny, critical signals of $\Delta W_{\text{merged}}$ are rounded to zero by the discrete quantization grid, effectively ``silencing'' the adapter's contribution entirely. 

To deconstruct this phenomenon and restore the promise of deployable model merging, we adopt the perspective of a critical methodologist. We refuse to accept SOTA model-merging claims that rely on high-precision evaluation strawmen. Instead, we propose a rigorous, multi-axial auditing and mitigation framework:

\begin{itemize}
  \item \textbf{The Re-Quantization Auditing (RQA) Framework:} We establish a systematic, multi-axial audit of model merging under quantization constraints. We evaluate performance across multiple quantization granularities (Per-Tensor vs. Per-Channel), bit-widths (4-bit vs. 8-bit), and formats (Symmetric vs. Asymmetric).
  \item \textbf{Discovery of the Quantization Granularity Bifurcation:} Our multi-axial audit reveals a major, previously uncharacterized bifurcation. While naive re-quantization is nearly lossless under standard per-channel grids (dropping only 0.15\% to 0.30\% accuracy), it collapses catastrophically under aggressive per-tensor constraints (dropping up to 8.6\% accuracy). We deconstruct the row-wise variance of task vectors to explain why per-channel grids provide inherent protection, exposing the per-tensor silence as a highly localized quantization artifact.
  \item \textbf{Deconstruction of the Representation Scale Preservation Dilemma (SAWS):} We evaluate Scale-Adaptive Weight Shifting (SAWS), a data-free closed-form scaling method. Crucially, we expose a fundamental mathematical contradiction in naive scale preservation: attempting to scale back the entire layer output to preserve activation scale collapses the pre-trained base model representations. We clarify that SAWS succeeds via ``selective task-vector boosting'' rather than true activation-scale preservation, deconstructing this architectural trade-off.
  \item \textbf{Exposing the Robustness and Limits of Test-Time Optimization (QA-ACS and AdaMerging):} We evaluate optimization-based test-time adaptation methods like AdaMerging and our proposed Quantization-Aware Adapter Coefficient Search (QA-ACS) to optimize merging coefficients directly through the quantization operator using the Straight-Through Estimator (STE). Our audit reveals that, contrary to full-precision assumptions, optimization-based approaches are highly robust, achieving the best overall performance under standard per-channel configurations. We critically analyze how unconstrained prediction entropy minimization can exhibit slight local instability (minor entropy decay) under severe per-tensor discretization noise, and demonstrate how this can be completely stabilized via basic coefficient regularization or supervised tuning.
\end{itemize}

We evaluate our framework on a Vision Transformer backbone across four image classification datasets (MNIST, FashionMNIST, CIFAR-10, SVHN). By shedding light on the re-quantization silence and deconstructing the failures of both data-free and optimization-based mitigations, this work provides a transparent and self-critical path forward for deployable, multi-task deep learning.